\section{Four campaign case studies}
\label{sec:cases}

\subsection{F1: an omitted gradient channel}
\label{sec:f1}

The decoy-gradient protocol mixes real activation rows with decoys. The
trusted client sends protected activations to the cloud; the cloud returns
its output; and the client, which holds the private loss, returns the output
gradient to the cloud. Because the loss ignores decoys, their gradient rows
are exactly zero. The resulting support pattern identifies which rows bear
the loss. This is a metadata disclosure of the real/decoy partition, not
token-identity reconstruction by zeros alone\citep{paper1} (\artifact{E1}).

The source reports exact partition recovery on 4,096 of 4,096 inspected
frames \emph{per run} across nine seeds. Its frequent-token content probe
has a separate effect of $+0.65$--$1.50$ percentage points over a
constant-guess baseline. These quantities have different denominators and
cannot be merged into a single recovery claim. In six utility-qualified
gradient-open cells, comprising three seeds on each of two datasets,
forward-only privacy and utility gates pass while the joint view exceeds
the pre-set $+1.0$ percentage-point privacy gate (\artifact{E1}).

The forward-only non-detection is therefore out-of-scope for the cloud's
joint-view privacy claim. The audit preserves the forward result and the
newly covered breach as separate findings.
The same source reports that per-row clipping with Gaussian noise removes
the scored partition disclosure: 0 of 1,024 inspected frames disclose the
partition, with agreement 0.400 at a 32/80 base rate. Approximately
$0.01$ nats is the associated held-out cross-entropy \emph{utility cost},
not residual leakage. This is scoped mitigation evidence under the scored
probes; attacks accumulating observations across training and other omitted
families remain untested. The case supports checking the joint transcript
and contradicts a universal dismissal of perturbation-based mitigation.

\subsection{F2: a reconstruction claim scored as prediction}
\label{sec:f2}

The boundary-replacement audit found that a representation computed from
\texttt{ids[:, :-1]} was scored against labels from \texttt{ids[:, 1:]}.
Thus a row representing $x_t$ was evaluated against $x_{t+1}$. The recorded
attack scores measure next-token predictive information, whereas the
headline privacy claim concerned reconstruction of the represented input.
A high positive-control score under this alignment cannot validate the
different reconstruction target. The current-token privacy conclusion is
invalid; the existing scores may be retained as next-token probe results
(\artifact{B1}).

The earlier independent artifact review also reports that separately
trained zero-return controls match or outperform the trained remote model
in all six clean corpus--seed comparisons. We retain this as a
\emph{source-reported} result, not a new recomputation. It fails to
demonstrate that learned remote computation is necessary for the observed
utility endpoints. It does not show that removing a return from an
already-trained model has literally no effect; that intervention differs
from training a matched model without the return (\artifact{B1}).

A further metric mismatch reinforces the need to name the estimand. In
one matched checkpoint evaluation, NLL rises by approximately $0.034$
nats and passes its utility rule, while teacher-forced answer-token top-1
match falls from $34.79\%$ to $18.15\%$, a $16.64$ percentage-point drop.
This is a specific answer-token metric, not standard multiple-choice task
accuracy or a replicated deployed-protocol result. The audit narrows the
quality and delegation claims without inventing corrected reconstruction
results that the historical evaluator did not produce (\artifact{B1}).

\subsection{F3: incomplete qualification and genuine rejections}
\label{sec:f3}

The trilemma campaign completed 232 exploratory runs and rejected all
19 configurations under its frozen exploratory rules. Reconciliation of
the cell table yields 13 qualified privacy breaches, five utility-gate
failures, and seven learning-gain failures. These reasons overlap. Six
configurations have undetermined privacy but are independently rejected
on utility or gain; width-64 is rejected on gain alone (\artifact{R1}).

Only four of 24 required instruments qualified: input-byte probes on
inbound and outbound gradients in the two recorded defence categories.
This missing coverage prevents a joint scoped privacy pass. It does not
invalidate the breaches detected by applicable qualified probes, nor the
independent utility or gain failures. The appropriate outcome is a finite
empty feasible set under the recorded rules, with incomplete privacy
coverage, rather than an explanation that all rejections arose from faulty
instruments (\artifact{R1}).

The span-6 cell illustrates why the inequality and rejection reasons must
be checked together. Its four-seed mean utility cost is
$0.0999705344$ nats and its upper 95\% Student-$t$ limit is
$0.1195024415$. Both \emph{pass} the respective $0.10$ and $0.12$
thresholds. The small $2.95\times10^{-5}$-nat mean margin lies on the
passing side. Its rejection reasons are gain failure and qualified privacy
breach, not utility failure. Four-seed intervals remain descriptive and
unadjusted; they do not establish population-level impossibility
(\artifact{R1}).

The subsequent eight-seed confirmation packet was proposed, not frozen,
admitted, or executed. Its $0.263772$ all-family power upper bound is a
conditional planning calculation using exploratory means and standard
deviations. It is not a probability that the 19 rejections are false.
Its sensitivity to variance assumptions and its explicit planning status
prevent its promotion into confirmatory outcome evidence (\artifact{R1}).

\subsection{F4: a calibrated finite negative result}
\label{sec:f4}

The defence-redesign pilot completed a frozen, exploratory 66-run campaign
covering seven cells with three paired seeds per cell and 16,384-window
panels. No cell was pilot-eligible. All seven pass the remote-gain rule;
five pass utility, while d3-q1 and d3-q2 fail it. The recorded composite
privacy verdicts are four failed and three undetermined. We preserve this
binding-based distinction rather than relabelling every missing
qualification as either success or breach (\artifact{P4}).

The pilot nevertheless supplies a substantive negative finding. Reported
worst-seed punctuation-property advantages are $0.79$--$0.89$ where
qualified, against the $0.02$ acceptance threshold; next-byte-class
advantages are $0.097$--$0.112$. Punctuation uses correlation advantage,
whereas next-byte structure uses classification-accuracy advantage.
Here next-byte prediction is an explicitly declared target, unlike F2's
use of that target to claim current-token reconstruction resistance.
A large detected signal relative to the policy gate is evidence against
the tested defence setting; its ratio alone does not demonstrate a broken
instrument or an impossible gate (\artifact{P4}).

Calibration constants have separate roles: the property positive-control
floor is $0.10$, the structure floor is $0.02$, the null band is $0.02$,
and privacy acceptance is $0.02$. A positive-control floor is not a
calibrated upper limit on undetected leakage. The qualification record
describes marginal null excursions as reliability evidence, not instrument
breakage. We therefore retain the finite negative result for the tested
mechanisms, strengths, budgets, and qualified targets, while withholding
claims about untested attacks, general token privacy, or universal
impossibility (\artifact{P4}).
